\documentclass[12pt,a4paper]{ctexart}

% ==== 字体与编码 ====
\usepackage{fontspec}
\setmainfont{Times New Roman}
\setsansfont{Arial}
\setmonofont{Consolas}
\setCJKmainfont{SimSun}[AutoFakeBold=2.5]
\setCJKsansfont{SimHei}
\setCJKmonofont{KaiTi}

% ==== 数学 ====
\usepackage{amsmath,amssymb,amsthm}
\usepackage{mathtools}

% ==== 页面布局 ====
\usepackage[top=2.5cm,bottom=2.5cm,left=2.8cm,right=2.8cm]{geometry}

% ==== 超链接 ====
\usepackage[colorlinks=true,linkcolor=blue,citecolor=blue,urlcolor=blue]{hyperref}

% ==== 交换图 ====
\usepackage{tikz-cd}

% ==== 列表 ====
\usepackage{enumitem}

% ==== 表格 ====
\usepackage{booktabs,array}

% ==== 定理环境 ====
\theoremstyle{definition}
\newtheorem{definition}{定义}[section]
\newtheorem{example}[definition]{例}
\newtheorem{remark}[definition]{注记}
\newtheorem{theorem}[definition]{定理}

% ==== 自定义命令 ====
\newcommand{\cat}[1]{\mathbf{#1}}
\newcommand{\Set}{\cat{Set}}
\newcommand{\id}{\mathrm{id}}
\newcommand{\ob}[1]{\mathrm{Ob}(#1)}
\newcommand{\Hom}{\mathrm{Hom}}

\title{\textbf{伴随函子的严格定义：自然同构视角}\\[0.3em]
\large ——从直观的"平滑"到严格的自然同构}
\author{范畴论学习笔记 · 范畴专栏}
\date{}

\begin{document}
\maketitle

\tableofcontents
\newpage

% ==================================================================
\section{引言：为何拒绝"平滑"而选择"自然同构"？}
% ==================================================================

在初学伴随函子时，常听到这样的直观描述：存在一族"平滑"的双射，在 $X$ 和 $Y$ 之间"自然变化"。然而数学的严谨性容不得半点含糊——"平滑"仅仅是一个为了直观而妥协的词汇。在范畴论的严格语境下，它必须被精确定义为\textbf{自然同构（Natural Isomorphism）}。

本文的目标是：将伴随函子基于 Hom-集的核心命题，逐层拆解为严谨的代数结构，直至每一个交换图和每一条等式都精确可验证。

\newpage

% ==================================================================
\section{核心命题：Hom-集双射的自然同构}
% ==================================================================

\begin{definition}[伴随函子，Hom-集版本]\label{def:adjunction-hom}
给定范畴 $\cat{C}$ 和 $\cat{D}$，以及函子 $F: \cat{C} \to \cat{D}$ 和 $G: \cat{D} \to \cat{C}$。

$F \dashv G$（称 $F$ 是 $G$ 的\textbf{左伴随}，$G$ 是 $F$ 的\textbf{右伴随}），当且仅当存在一族双射：
\begin{equation}\label{eq:phi-family}
\Phi_{X,Y}: \Hom_{\cat{D}}(F(X), Y) \xrightarrow{\; \sim \;} \Hom_{\cat{C}}(X, G(Y))
\end{equation}
并且这族双射在 $X \in \ob{\cat{C}}$ 和 $Y \in \ob{\cat{D}}$ 上是\textbf{自然的（Natural）}。
\end{definition}

\textbf{关键追问：}"在 $X$ 和 $Y$ 上是自然的"究竟意味着什么？仅仅说"自然"是一个未经定义的模糊词汇，下面给出严格拆解。

\newpage

% ==================================================================
\section{严谨拆解：双函子与自然性}
% ==================================================================

\subsection{两侧定义的是双函子}

要谈论"自然变换"，必须先明确它连接的是哪两个函子。等式 \eqref{eq:phi-family} 的两边，实际上定义了两个从乘积范畴 $\cat{C}^{\mathrm{op}} \times \cat{D}$ 到集合范畴 $\Set$ 的\textbf{双函子（Bifunctors）}：

\begin{enumerate}[label=(\arabic*),nosep]
    \item \textbf{左侧双函子} $L(-, -) = \Hom_{\cat{D}}(F(-), -)$
    \item \textbf{右侧双函子} $R(-, -) = \Hom_{\cat{C}}(-, G(-))$
\end{enumerate}

注意第一变量是逆变的（属于反范畴 $\cat{C}^{\mathrm{op}}$），第二变量是协变的。

\begin{remark}
双函子的逆变/协变区分至关重要：
\begin{itemize}[nosep]
    \item 在 $L$ 中：$F$ 作用于第一变量（协变），但 $F$ 位于 Hom-集的第一个位置，而 Hom-集的第一位置本身是逆变的。因此 $L$ 在第一变量上整体逆变。
    \item 在 $R$ 中：第一变量直接位于 Hom-集的第一个位置，因此逆变；第二变量 $G(-)$ 位于 Hom-集的第二个位置（协变）。
\end{itemize}
所以两个双函子都具有类型 $\cat{C}^{\mathrm{op}} \times \cat{D} \to \Set$，这才能谈论它们之间的自然变换。
\end{remark}

\subsection{自然性的精确定义}

"$\Phi$ 是自然的"这一表述，其严谨等价命题是：

\begin{center}
\framebox{\textbf{$\Phi$ 是双函子 $L$ 到 $R$ 的自然同构。}}
\end{center}

这意味着 $\Phi$ 是函子范畴 $[\cat{C}^{\mathrm{op}} \times \cat{D}, \Set]$ 中的一个可逆态射。

落实到具体的态射和交换图上，这要求 $\Phi$ 必须同时满足\textbf{对 $X$ 的自然性（逆变）}和\textbf{对 $Y$ 的自然性（协变）}。

\newpage

\subsection{对变量 $X$ 的自然性（逆变方向）}

任取 $\cat{C}$ 中的态射 $f: X' \to X$，以及任意对象 $Y \in \cat{D}$ 和态射 $h: F(X) \to Y$。

考虑两条从 $\Hom_{\cat{D}}(F(X), Y)$ 到 $\Hom_{\cat{C}}(X', G(Y))$ 的路径：

\begin{itemize}[nosep]
    \item \textbf{路径一（先复合在 $\cat{D}$，再双射）：}
    先在 $\cat{D}$ 中将 $h$ 与 $F(f)$ 前复合，得到 $h \circ F(f): F(X') \to Y$，然后应用 $\Phi_{X',Y}$。
    \[
    h \;\xmapsto{-\circ F(f)}\; h \circ F(f) \;\xmapsto{\Phi_{X',Y}}\; \Phi_{X',Y}(h \circ F(f))
    \]

    \item \textbf{路径二（先双射，再复合在 $\cat{C}$）：}
    先应用 $\Phi_{X,Y}$ 得到 $\Phi_{X,Y}(h): X \to G(Y)$，然后在 $\cat{C}$ 中将其与 $f$ 后复合。
    \[
    h \;\xmapsto{\Phi_{X,Y}}\; \Phi_{X,Y}(h) \;\xmapsto{-\circ f}\; \Phi_{X,Y}(h) \circ f
    \]
\end{itemize}

自然性要求这两条路径结果相等，从而得到\textbf{等式一}（对 $X$ 的自然性）：
\begin{equation}\label{eq:natural-X}
\boxed{\; \Phi_{X',Y}\bigl(h \circ F(f)\bigr) = \Phi_{X,Y}(h) \circ f \;}
\end{equation}

用交换图表示：
\[
\begin{tikzcd}[column sep=4em, row sep=3em]
\Hom_{\cat{D}}(F(X), Y) \arrow[r,"\Phi_{X,Y}","\sim"'] \arrow[d,"-\circ F(f)"'] & \Hom_{\cat{C}}(X, G(Y)) \arrow[d,"-\circ f"] \\
\Hom_{\cat{D}}(F(X'), Y) \arrow[r,"\Phi_{X',Y}","\sim"'] & \Hom_{\cat{C}}(X', G(Y))
\end{tikzcd}
\]

\subsection{对变量 $Y$ 的自然性（协变方向）}

任取 $\cat{D}$ 中的态射 $g: Y \to Y'$，以及任意对象 $X \in \cat{C}$ 和态射 $k: F(X) \to Y$。

考虑两条从 $\Hom_{\cat{D}}(F(X), Y)$ 到 $\Hom_{\cat{C}}(X, G(Y'))$ 的路径：

\begin{itemize}[nosep]
    \item \textbf{路径一（先后复合在 $\cat{D}$，再双射）：}
    先在 $\cat{D}$ 中将 $g$ 与 $k$ 后复合，得到 $g \circ k: F(X) \to Y'$，然后应用 $\Phi_{X,Y'}$。
    \[
    k \;\xmapsto{g\circ-}\; g \circ k \;\xmapsto{\Phi_{X,Y'}}\; \Phi_{X,Y'}(g \circ k)
    \]

    \item \textbf{路径二（先双射，再复合在 $\cat{C}$）：}
    先应用 $\Phi_{X,Y}$ 得到 $\Phi_{X,Y}(k): X \to G(Y)$，然后与 $G(g)$ 后复合。
    \[
    k \;\xmapsto{\Phi_{X,Y}}\; \Phi_{X,Y}(k) \;\xmapsto{G(g)\circ-}\; G(g) \circ \Phi_{X,Y}(k)
    \]
\end{itemize}

自然性要求这两条路径结果相等，从而得到\textbf{等式二}（对 $Y$ 的自然性）：
\begin{equation}\label{eq:natural-Y}
\boxed{\; \Phi_{X,Y'}\bigl(g \circ k\bigr) = G(g) \circ \Phi_{X,Y}(k) \;}
\end{equation}

用交换图表示：
\[
\begin{tikzcd}[column sep=4em, row sep=3em]
\Hom_{\cat{D}}(F(X), Y) \arrow[r,"\Phi_{X,Y}","\sim"'] \arrow[d,"g\circ-"'] & \Hom_{\cat{C}}(X, G(Y)) \arrow[d,"G(g)\circ-"] \\
\Hom_{\cat{D}}(F(X), Y') \arrow[r,"\Phi_{X,Y'}","\sim"'] & \Hom_{\cat{C}}(X, G(Y'))
\end{tikzcd}
\]

\begin{remark}[两条自然性条件的独立性]
等式一和等式二是相互独立的：前者约束 $X$ 方向的前复合行为，后者约束 $Y$ 方向的后复合行为。一个双射族 $\{\Phi_{X,Y}\}$ 要成为自然同构，必须\textbf{同时}满足这两条等式。只满足其中一条不足以保证自然性。
\end{remark}

\newpage

% ==================================================================
\section{完整严格定义与统一方程}
% ==================================================================

\subsection{伴随函子的完整定义}

综合上述所有条件，给出完全的严格定义：

\begin{definition}[伴随函子的严格定义]\label{def:adjunction-full}
函子 $F: \cat{C} \to \cat{D}$ 是函子 $G: \cat{D} \to \cat{C}$ 的\textbf{左伴随}，当且仅当存在一个自然同构：
\begin{equation}
\Phi: \Hom_{\cat{D}}(F(-), -) \;\cong\; \Hom_{\cat{C}}(-, G(-))
\end{equation}
作为函子范畴 $[\cat{C}^{\mathrm{op}} \times \cat{D}, \Set]$ 中的同构。
\end{definition}

\subsection{统一方程（合并两条自然性）}

两条自然性条件可以合并为一条统一的方程。对于任意对象 $X, X' \in \cat{C}$，任意对象 $Y, Y' \in \cat{D}$，以及任意态射 $f: X' \to X$（在 $\cat{C}$ 中）、$g: Y \to Y'$（在 $\cat{D}$ 中）和 $h: F(X) \to Y$（在 $\cat{D}$ 中），下述方程成立：

\begin{equation}\label{eq:combined}
\boxed{\; \Phi_{X',Y'}\!\bigl(g \circ h \circ F(f)\bigr) = G(g) \circ \Phi_{X,Y}(h) \circ f \;}
\end{equation}

\textbf{推导：} 先对 $g$ 应用等式二，再对 $f$ 应用等式一：
\begin{align*}
\Phi_{X',Y'}(g \circ h \circ F(f))
&= G(g) \circ \Phi_{X',Y}(h \circ F(f)) \tag{等式二}\\
&= G(g) \circ \bigl(\Phi_{X,Y}(h) \circ f\bigr) \tag{等式一}\\
&= G(g) \circ \Phi_{X,Y}(h) \circ f.
\end{align*}

用交换图（3D 视角的"管状图"）表示：
{\small\[
\begin{tikzcd}[column sep=5em, row sep=3em]
\Hom_{\cat{D}}(F(X), Y) \arrow[r,"\Phi_{X,Y}","\sim"'] \arrow[d,"-\circ F(f)"'] \arrow[dd, bend right=70, "g\circ-"'] &
\Hom_{\cat{C}}(X, G(Y)) \arrow[d,"-\circ f"] \arrow[dd, bend left=70, "G(g)\circ-"] \\
\Hom_{\cat{D}}(F(X'), Y) \arrow[r,"\Phi_{X',Y}","\sim"'] \arrow[d,"g\circ-"'] &
\Hom_{\cat{C}}(X', G(Y)) \arrow[d,"G(g)\circ-"] \\
\Hom_{\cat{D}}(F(X'), Y') \arrow[r,"\Phi_{X',Y'}","\sim"'] &
\Hom_{\cat{C}}(X', G(Y'))
\end{tikzcd}
\]}
此图的两个正方形面分别对应等式一（前面）和等式二（右面），而外围的矩形面正是统一方程 \eqref{eq:combined}。

\subsection{逐条条件对照}

\begin{table}[h]
\centering
\caption{伴随函子定义的条件摘要}
\begin{tabular}{@{}l|l|l@{}}
\toprule
\textbf{条件} & \textbf{方程} & \textbf{含义} \\
\midrule
双射性 & $\Phi_{X,Y}$ 是一一映射 & 态射集之间存在一一对应 \\
对 $X$ 的自然性 & $\Phi_{X',Y}(h \circ F(f)) = \Phi_{X,Y}(h) \circ f$ & 双射与 $\cat{C}$ 中的前复合协调 \\
对 $Y$ 的自然性 & $\Phi_{X,Y'}(g \circ k) = G(g) \circ \Phi_{X,Y}(k)$ & 双射与 $\cat{D}$ 中的后复合协调 \\
统一方程 & $\Phi_{X',Y'}(g \circ h \circ F(f)) = G(g) \circ \Phi_{X,Y}(h) \circ f$ & 上述两条的合并形式 \\
\bottomrule
\end{tabular}
\end{table}

\newpage

% ==================================================================
\section{总结与评注}
% ==================================================================

\subsection{核心要点}

\begin{enumerate}[nosep]
    \item \textbf{自然同构}连接的是两个\textbf{双函子} $\Hom_{\cat{D}}(F(-), -)$ 和 $\Hom_{\cat{C}}(-, G(-))$，它们都是从 $\cat{C}^{\mathrm{op}} \times \cat{D}$ 到 $\Set$ 的函子。
    \item \textbf{自然性}拆解为两个独立的方向：对 $X$ 的逆变自然性和对 $Y$ 的协变自然性。
    \item \textbf{统一方程} $\Phi_{X',Y'}(g \circ h \circ F(f)) = G(g) \circ \Phi_{X,Y}(h) \circ f$ 囊括了全部自然性条件，是伴随函子最简洁的等价定义。
    \item 所有映射在态射的前复合与后复合操作下都完全兼容（Commutative），不存在任何"近似"或"平滑"——每一步都是精确的代数等式。
\end{enumerate}

\subsection{为何这条定义比"单位-余单位"定义更基础？}

基于 Hom-集的伴随定义（定义 \ref{def:adjunction-full}）是范畴论中最原始、最通用的定义。而基于单位 $\eta$ 和余单位 $\varepsilon$ 的伴随定义是它的等价形式：

\begin{itemize}[nosep]
    \item 单位：$\eta_X = \Phi_{X,F(X)}(\id_{F(X)}): X \to G(F(X))$；
    \item 余单位：$\varepsilon_Y = \Phi_{G(Y),Y}^{-1}(\id_{G(Y)}): F(G(Y)) \to Y$。
\end{itemize}

Hom-集版本的优势在于：它直接暴露了伴随的本质——态射集之间的自然双射——不需要额外的三角形恒等式推导，从而在教学中作为入门定义更为清晰。

\subsection{与后续学习笔记的连接}

本文是范畴专栏中"基础构件"系列的一部分。掌握了伴随函子的严格定义后，下一讲将进入：
\begin{itemize}[nosep]
    \item 伴随函子的等价刻画（单位-余单位形式及其三角形恒等式）；
    \item 伴随函子的例子（自由-遗忘伴随、积-指数伴随等）；
    \item 单子（Monad）作为伴随的导出结构。
\end{itemize}

\vspace{1em}
\begin{center}
\rule{0.5\textwidth}{0.5pt}\\[0.5em]
\textbf{本节完}
\end{center}

\end{document}